\documentclass{article}
\usepackage[margin=1in]{geometry}
\usepackage{enumitem}
\usepackage{hyperref}

\title{SolFilOsc Project Reorganization Report}
\date{2026-05-12}

\begin{document}
\maketitle

\section{Purpose}
The project was reorganized so the preprocessing code and the CNN-based
oscillation analysis live in importable Python modules under \texttt{src/}.
The numerical workflow was preserved; the change is structural and
documentational.

\section{Code Layout Changes}
\begin{itemize}[leftmargin=*]
  \item Moved the former \texttt{data\_processing/} scripts to
  \texttt{src/solfilosc/data\_processing/}.
  \item Moved the mapping utilities to \texttt{src/solfilosc/mapping/}.
  \item Split the notebook implementation into focused modules under
  \texttt{src/solfilosc/analysis/}: constants, CNN/noise-model utilities,
  conformal-prediction calibration, degradation helpers, ROI helpers,
  detection/component extraction, event clustering, plotting, writers, and
  the pipeline driver.
  \item Replaced \texttt{notebooks/analysis\_with\_CNN.ipynb} with a thin
  interactive launcher that imports the packaged analysis pipeline.
  \item Converted the unused legacy \texttt{fits\_obj.py} IDL-style reference
  into an explicit Python placeholder because the previous file was not valid
  Python and was not imported by the active pipeline.
\end{itemize}

\section{Import and Execution Changes}
\begin{itemize}[leftmargin=*]
  \item Updated internal imports from folder-relative assumptions to package
  imports such as \texttt{solfilosc.mapping} and relative imports inside
  \texttt{solfilosc.data\_processing}.
  \item Updated \texttt{preprocessing.sh} to run preprocessing stages with
  \texttt{python3 -m solfilosc.data\_processing.<module>} and to add
  \texttt{src/} to \texttt{PYTHONPATH}.
  \item Added \texttt{pyproject.toml} so the package can be installed in
  editable mode with \texttt{pip install -e .}. The analysis command is
  exposed as \texttt{solfilosc-analysis} through a lightweight CLI wrapper.
  \item Updated \texttt{preprocess\_data.py} so the default raw FITS location
  is project-local: \texttt{data/raw/<month>/<day>/}. A custom raw-data
  directory can still be supplied as an optional fifth command-line argument.
  If the expected folder or \texttt{*.fits.fz} files are missing, the script
  now raises an explicit setup error instead of failing later by indexing an
  empty file list.
\end{itemize}

\section{Documentation Changes}
\begin{itemize}[leftmargin=*]
  \item Rewrote \texttt{README.md} to document the repository layout,
  installation procedure, preprocessing pipeline, filament segmentation
  dependency, CNN analysis execution, and operational caveats.
  \item Expanded the preprocessing documentation with the exact positional
  command-line arguments, default raw-data layout, optional custom raw-data
  path, FITS filename/header assumptions, and FilamentSeg citation details.
\end{itemize}

\section{Functionality}
No intentional numerical or scientific-method changes were made. The
preprocessing stages still perform file filtering, limb-darkening/background
correction, derotation/data-cube creation, and post-processing. The analysis
still performs CNN noise-model inference, conformal-prediction thresholding,
multiscale period detection, component extraction, event clustering, and
plot/table generation.

\end{document}
